\section{Módulo de Cuotas}
Este módulo está diseñado para llevar un seguimiento preciso de las compras realizadas a plazo. El enfoque del modelo de datos prioriza el flujo de caja recurrente sobre la deuda histórica: 	extbf{el sistema no almacena el precio original de la compra}, sino únicamente el número de cuotas y el valor de cada una.

Las cifras agregadas que muestra la interfaz (total a pagar, abonado y deuda pendiente) se 	extbf{derivan} de esos dos datos como 	extit{cuotas} $	imes$ 	extit{valor}. Conviene notar que esa multiplicación no equivale al precio de lista del producto, ya que incorpora los intereses del financiamiento: representa el desembolso total comprometido, que es la magnitud relevante para el flujo de caja.

\subsection{Proyección de la Carga Mensual}
El módulo incorpora una proyección de lo que el usuario pagará en cuotas durante los próximos doce meses. Dado que cada producto aporta su cuota únicamente mientras le resten pagos, la serie solo puede descender: \textbf{cada escalón hacia abajo corresponde a una deuda que termina}.

Esta lectura responde la pregunta central del módulo, que el mero listado de productos no resuelve: cuándo se libera capacidad de pago y en qué cuantía. Los indicadores que acompañan al gráfico precisan la \textbf{carga mensual} vigente, la \textbf{deuda pendiente} total, la \textbf{próxima liberación} (mes y monto en que disminuye la carga) y la fecha en que el usuario queda \textbf{libre de cuotas}.

\subsection{Registro de Elementos}
Para dar de alta un nuevo pago en cuotas, el usuario debe registrar la siguiente información del producto:
\begin{itemize}
    \item \textbf{Producto:} Nombre o descripción de la compra.
    \item \textbf{Tienda:} Comercio o entidad donde se adquirió.
    \item \textbf{Cantidad de cuotas:} Número total de pagos acordados.
    \item \textbf{Valor de la cuota:} Monto recurrente a pagar en cada periodo.
    \item \textbf{Medio de pago:} Instrumento con el que se paga el compromiso. Es opcional y se declara una sola vez, puesto que no cambia de una cuota a otra; los gastos que el sistema genera al marcar cada pago lo \textbf{heredan} automáticamente.
\end{itemize}

\subsection{Gestión e Interacción}
La interfaz permite interactuar con cada registro listado. Al presionar sobre un elemento, se despliega un menú de acciones para gestionar el estado de la deuda:
\begin{itemize}
    \item \textbf{Marcar cuota:} Acción rápida que registra un pago efectuado, incrementando el contador de cuotas pagadas y actualizando el progreso.
    \item \textbf{Editar:} Permite la modificación integral del registro. El usuario puede actualizar el producto, la tienda, el total de cuotas, el valor de las mismas y el \textbf{número de cuotas pagadas}. Esta última propiedad es crucial para brindar flexibilidad, permitiendo revertir acciones en caso de presionar "Marcar cuota" por accidente.
    Al reducir el contador de cuotas pagadas, el sistema \textbf{elimina también los gastos que esos marcados habían generado} en el libro mayor, de modo que la corrección no deje movimientos fantasma en el historial. El formulario advierte de ello antes de confirmar. El módulo ofrece además un atajo de \textit{Deshacer última cuota} para el caso más frecuente, que evita abrir el formulario completo.
    \item \textbf{Eliminar:} Borra el registro de la base de datos de forma permanente. Los gastos ya generados por el producto \textbf{se conservan} en el historial, puesto que corresponden a dinero efectivamente desembolsado; únicamente pierden la trazabilidad hacia el producto de origen.
\end{itemize}